\documentclass{article}
\usepackage{graphicx}
\usepackage{amsmath}
\usepackage{amssymb}
\usepackage{url}
\usepackage{booktabs}
\usepackage[margin=1in]{geometry}
\usepackage{natbib}

\title{Captive Cross-Border Demand and Rural Agriculture:\\
Cannabis Revenue, Farmland Values, and Crop Substitution\\
in Border Counties of Legal-Market States}
\author{Mahyar Rezazadeh \\ Babak Sahragard}
\date{August 2026}

\begin{document}

\maketitle

\begin{abstract}
\noindent We test whether counties in legal-cannabis states that are close to a non-legal-state border --
and can therefore absorb captive demand from residents of the neighboring illegal-market state --
show different economic outcomes than interior counties of the same legal states. This is an
entry-game/border-competition design: both treatment and control counties are in already-legal states;
the comparison is border versus interior, not legal versus illegal. Using a 70-mile border threshold
across the 23 states with an operating recreational retail market as of August 2026, we test four
outcomes: cannabis retail revenue per capita and revenue per dispensary store (Colorado, Washington,
Oregon, and Montana -- the states with usable county-level revenue data), farmland asset value per acre
(USDA NASS Census of Agriculture, available nationally), and harvested cropland acreage (Oregon). Revenue
per capita, farmland value, and crop substitution mostly come back null or inconclusive in Colorado,
Washington, and Oregon. Revenue \emph{per store}, however, is significantly higher in border counties in
all three of those states -- pooling them strengthens rather than weakens the result -- robust to excluding
a single-county outlier in each state and to estimating each state separately, even though total
county-level revenue and dispensary counts show no border effect. An event study by state market age shows
this premium strengthens rather than fades as a legal market matures, and two further mechanism checks --
a continuous distance dose-response and, more robustly, a measure of the total non-legal population within
70 miles rather than distance to the nearest non-legal county alone -- both point in the direction the
cross-border demand hypothesis predicts, with the population-weighted measure surviving county-clustered
inference; the per-store result is, however, sensitive to widening the border threshold from 70 to 90
miles, a limitation we report directly. Montana reverses this pattern entirely: it is the one state in the
sample with no statewide cap on dispensary licenses, and its revenue-per-capita coefficient is the one
county-aggregate estimate in the paper that is positive and significant under county-clustered inference,
while its own revenue-per-store estimate is a precise null -- consistent with captive demand being absorbed
through new store entry rather than higher revenue per existing store when entry is not capped. A placebo
check against Colorado's general (non-cannabis) retail sales -- the single most important robustness check
in the paper -- finds no comparable border/interior gap in ordinary retail activity over the same period,
arguing against the alternative explanation that border counties simply transact more retail business of
any kind. A second placebo check, using alcohol sales in Oregon and Montana (a closer analogue to cannabis
than general retail, being similarly regulated and excise-taxed), is more mixed: Oregon reproduces the
general-retail null and, if anything, shows a lower per-store premium near the border, while Montana shows
a positive border/interior gap in alcohol revenue per capita that is significant under heteroskedasticity-robust
but not county-clustered inference -- we report this directly rather than treating it as a second clean
pass. We also report, as suggestive context rather than a formal test, a directly relevant historical
data point: Washington's 2014 initial license lottery shows one of this paper's own border counties
(Spokane) was sharply oversubscribed for cannabis retail licenses, consistent with real unmet entry demand
resolved by a binding cap rather than by market entry. We read the overall pattern as consistent with a
demand response that is real but absorbed by a binding constraint on store entry where such a constraint
exists (state licensing caps), exactly the explanation proposed when a companion analysis found no effect
of border proximity on dispensary density: captive demand does not appear to draw in more stores where
entry is capped, but it does appear to make each existing store busier -- and where entry is not capped, as
in Montana, it draws in more stores instead. Farmland value is null in the panel specification with state and year fixed
effects, though a continuous distance measure disagrees in sign with the binary threshold. An event study
by state market age, extended here to include the pre-launch years that NASS's Census waves make available
for several states, shows a significant \emph{negative} border/interior gap emerging after several years of
legal-market operation -- but also shows the pre-launch period itself is not flat, with significant swings
in both directions before legalization even occurs, which complicates treating the post-launch decline as a
confident causal estimate rather than a continuation of pre-existing volatility. Crop substitution
in Oregon is essentially untestable at this threshold: only three of Oregon's 36 counties fall within 70
miles of its one non-legal neighbor, and the state's actual cannabis cultivation hub sits over 300 miles
away, near the California border instead. Finally, we subject the revenue-per-store result to three checks
that move beyond a fixed cross-sectional comparison. Following \citet{HansenMillerWeber2020}, who show
Washington border-with-Oregon sales fell 36 percent when Oregon's own market opened, we test whether
Colorado counties bordering New Mexico lost their premium after New Mexico's 2022 retail market opened: they
did, by 38.6 percent, a magnitude close to \citeauthor{HansenMillerWeber2020}'s own estimate at an
independent border; our own lower-frequency replication of their original Washington-Oregon setting,
however, does not reproduce their result, which we attribute to this paper's coarser annual data rather
than a genuine absence of the mechanism there. We also find that dispensaries do not open faster in border
counties initially, but a second dispensary follows significantly faster once a border county has a first
one, and that dispensaries within border counties are located measurably closer to the actual state line
than the county's own geography would predict -- direct evidence that sellers respond to, rather than are
merely correlated with, cross-border demand.
\end{abstract}

\section{Introduction}

A recurring hypothesis in this research program is that counties in legal-cannabis states located near a
non-legal-state border can capture demand from residents of the neighboring state, who may travel some
distance to purchase legally rather than rely on an illegal supplier at home. This is distinct from every
other design used elsewhere in this program: rather than asking how a non-legal county is affected by
proximity to a legal market, here \emph{both} treatment and control counties are already inside a legal
state, and the comparison is between counties near the state's outer, non-legal-adjacent border and
counties in the interior. A companion analysis using this design and 2020--2025 dispensary data found no
evidence that dispensary entry timing or density responds to this kind of border proximity. This paper
extends the same treatment/control logic to three further outcomes -- retail revenue, farmland value, and
crop substitution -- that speak to different economic channels than store counts alone: revenue captures
demand intensity even where entry itself is capped by state licensing rules; farmland value captures
capitalized expectations about a county's economic trajectory; and crop substitution asks whether
cannabis cultivation displaces traditional agriculture disproportionately in these captive-demand
counties.

\citet{SchwabButsic2020} provide the most direct related evidence, finding that a 2016 Humboldt County,
California ordinance making specific parcels eligible for a cannabis cultivation permit raised those
parcels' sale price roughly 84 percent relative to ineligible parcels. That result identifies the value of
\emph{cultivation-permit eligibility} within one county using bespoke parcel-level data; it does not speak
to whether simple \emph{proximity to a captive cross-border demand pool} moves outcomes at the county
level across many states, which is the question this paper asks instead.

We test all three outcomes using the identical treatment and control definition, described precisely in
Section~\ref{sec:data}. Revenue per capita, farmland value, and crop substitution find no robust support
for the captive-demand hypothesis, though, as we discuss, the reasons differ meaningfully across outcomes,
and one specification (the land-value event study by market age) surfaces a result opposite in sign to the
hypothesis rather than a simple null. Revenue \emph{per store}, by contrast, is a robust, outlier-independent,
multi-state positive result (Section~\ref{sec:revenue_perstore}). Section~\ref{sec:switchoff} then subjects
that result to three additional checks that do not rely on the same fixed cross-sectional comparison:
whether the premium switches off when a bordering state legalizes its own retail market, following
\citet{HansenMillerWeber2020}; whether sellers entered border counties faster; and whether they located
measurably closer to the state line itself.

\section{Data}
\label{sec:data}

\subsection{Treatment and control: border versus interior counties of legal states}

We identify all states with an operating recreational cannabis retail market as of August 5, 2026, using
each state's first-retail-sale date. Twenty-three states qualify: Colorado and Washington (both legal
since 2014) through Delaware (legal since August 2025). For every county in these 23 states, we compute
the great-circle distance from the county's population centroid to the nearest county located in a state
\emph{without} an operating recreational retail market, using the same centroid data and distance
methodology used throughout this research program. This yields 991 counties across the 23 legal states.

\textbf{Treatment} counties are those within 70 miles of a non-legal-state border -- close enough, under
this design's motivating logic, to plausibly absorb demand from residents of the neighboring illegal-market
state who travel up to roughly that distance to buy legally. \textbf{Control} counties are the interior,
``middle-of-the-state'' counties of the same 23 legal states: more than 70 miles from any non-legal
border. Both groups are, by construction, counties in states where cannabis is already legal; this is not
a legal-versus-illegal comparison, and a county's own eligibility to host cultivation or retail activity
plays no role in the treatment definition -- only its distance to a captive-demand-generating border.

Border-county availability varies sharply by state. Five of the 23 states -- Arizona, California, Rhode
Island, and (in the two outcome-specific panels below, which have slightly different state coverage due to
data-matching) Connecticut and Alaska -- have \emph{no} county within 70 miles of a non-legal border;
California's closest county to any non-legal state is 222 miles away, and Alaska's is 769 miles away,
since neither is bounded by a state that remains illegal. A sixth state, Nevada, is a near-miss rather than
a clean exclusion: its closest county, Lincoln (79.7 miles from Utah), and Elko (90.8 miles from Idaho)
both sit just outside the 70-mile cutoff, reflecting Nevada's large, sparsely populated counties rather
than genuine geographic remoteness from a non-legal border -- unlike California and Alaska, Nevada does
directly border two non-legal states (Idaho and Utah), and the population centroids used throughout this
paper simply place its nearest cities just past the line. Section~\ref{sec:threshold} reports a 90-mile
robustness check motivated directly by this near-miss. These states can contribute control observations
but no treated ones at the 70-mile threshold, and are excluded from the robustness checks that restrict to
states with genuine border variation.

\subsection{Outcome 1: cannabis retail revenue, per capita and per store}
\label{sec:data_revenue}

We reuse the county-fiscal-year revenue panel constructed for a companion analysis in this project: monthly
county-level marijuana sales from Colorado's Department of Revenue (2014--2025) and Washington's Liquor
and Cannabis Board (2015--2025). We extend this panel with two additional states' own regulator data:
Oregon's OLCC publishes a county-by-month sales series through its open-data portal (2016--2026, all
product and customer types combined), and Montana's Department of Revenue publishes monthly per-county
sales as a series of static report pages (no bulk download exists; we retrieved 53 of 54 available months
directly). Illinois, Ohio, and Maine publish statewide totals only; Massachusetts's county-level data
exists only behind an inaccessible interactive dashboard; New York's data is a regional cumulative total,
not a county panel; and no other state's regulator data proved retrievable at county granularity within
the scope of this project. Colorado, Washington, Oregon, and Montana are therefore the four states this
paper's revenue analysis draws on. Our primary revenue outcome is log(revenue per capita), using each
county's population as the normalizing denominator.

We additionally construct log(revenue per dispensary store). For Colorado and Washington this uses a
geocoded dispensary-license dataset (759 records, license number, address, and opening date, assembled for
a companion analysis of dispensary entry) point-in-polygon assigned to counties, counting dispensaries open
as of each year's own reporting cutoff (December 31 for Colorado's calendar-year reporting, June 30 for
Washington's fiscal-year reporting). The same geocoded dataset covers Oregon nominally, but proved
unreliable specifically for Oregon: several records geocode to counties in other states entirely, and
Baker County -- one of Oregon's three border counties -- has zero matched dispensaries despite recording
real revenue in every year from 2017 to 2025. We instead use OLCC's own license registry, which tags
county natively; because that registry's date field reflects each license's most recent renewal rather
than its original opening date, it cannot support a historical panel, so Oregon's revenue-per-store check
is a single 2025 cross-section (2025 revenue against currently-active retailer counts) rather than a
multi-year panel. Montana's revenue-per-store check is built the same way, from the Department of
Revenue's current licensed-dispensary list matched to county via each dispensary's city (a 99\% match
rate against a crosswalk built from Montana's incorporated municipalities). County-years with zero open
stores are dropped throughout, since revenue per store is undefined. This outcome asks a different
question than revenue per capita: whether individual outlets transact more business in border counties,
independent of how many people live there or how many stores have opened to serve them.

\subsection{Placebo outcome: general retail sales}

To test whether the revenue-per-store premium is specific to cannabis or simply reflects border counties
transacting more retail business of any kind, we construct a placebo outcome from the Colorado Department
of Revenue's county-by-industry retail sales filings \citep{CDORRetailByCounty}, which report gross retail
sales and the number of filing retailers by county, month, and industry across all sectors of the state
economy. We use the ``Total'' (all-industry) category, aggregated to complete calendar years,
2016--2025. Colorado's Department of Revenue discloses this level of detail only for the state's 11
largest counties by filing volume, which includes 3 of the 6 border counties in the main analysis
(Larimer, Mesa, and Weld) but not the smaller ones (Montezuma, Morgan, and Sedgwick); we note this
coverage gap explicitly rather than treating the placebo sample as identical to the main analysis's
border group.

\subsection{Outcome 2: farmland asset value per acre}

Farmland value is \emph{AG LAND, INCL BUILDINGS -- ASSET VALUE, MEASURED IN \$/ACRE} from the USDA NASS
Census of Agriculture, extracted at the county level for the 2012, 2017, and 2022 Census waves -- the only
years this variable is published below the state level. Because NASS's Census of Agriculture is a
national program, this outcome is available for essentially every county regardless of state, unlike
revenue: we match 947 of the 991 legal-state counties to complete land-value records in all three years,
across 21 of the 23 states.

\subsection{Outcome 3: harvested cropland acreage}

For crop substitution we use \emph{AG LAND, CROPLAND, HARVESTED -- ACRES}, again from the NASS Census of
Agriculture, for all 36 Oregon counties in 2012, 2017, and 2022. We focus on Oregon because it is the only
state for which we were able to independently verify a usable cannabis cultivation dataset (see the
companion analysis in this project for why); as Section~\ref{sec:results} shows, however, Oregon turns out
to have very little genuine border-county variation under the 70-mile threshold, since it borders only
one non-legal state (Idaho) at a considerable remove from most of its population.

\section{Empirical Strategy}
\label{sec:methods}

For revenue (per capita and per store) and farmland value, all three observed as county-year or
county-fiscal-year panels, we estimate
\[
y_{ist} = \alpha + \beta \, \text{Border}_{i} + \gamma \ln(\text{population}_{i}) + \delta \ln(\text{land area}_i)
+ \theta_t + \mu_s + \varepsilon_{ist},
\]
where $y_{ist}$ is log(revenue per capita), log(revenue per store), or log(land value per acre) for county
$i$ in state $s$ at time $t$, $\text{Border}_i \in \{0,1\}$ is the 70-mile treatment indicator (time-invariant, since a county's
distance to a non-legal border does not change), and $\theta_t$, $\mu_s$ are year and state fixed effects.
\textbf{Standard errors are clustered at the county level} throughout the panel specifications, since each
county contributes multiple correlated observations over time. We report both heteroskedasticity-robust
(HC1) and county-clustered standard errors for transparency; the clustered version is our preferred
inference. We also estimate a continuous version replacing $\text{Border}_i$ with $\ln(\text{distance to
nearest non-legal border}_i)$.

For revenue per store specifically, we additionally construct a market-size measure intended to test the
underlying mechanism more directly than distance to the nearest non-legal county alone: for each Colorado
and Washington county we sum the population of every non-legal-state county within 70 miles (not only the
nearest one), using a radius query against the national county centroid and population-estimate files
already used to build the treatment panel. This $\text{CaptivePop}_{70,i}$ measure distinguishes a county
adjacent to one small non-legal county from a county adjacent to a large non-legal metropolitan area, which
a nearest-county distance measure cannot; we replace $\text{Border}_i$ with $\ln(\text{CaptivePop}_{70,i}+1)$
and re-estimate the same specification.

The placebo outcome (general retail sales per capita) uses the identical specification restricted to
Colorado only, so the state fixed effect $\mu_s$ drops out; year fixed effects and the log(population)
control are retained.

Because border status does not switch over time -- a county cannot move from interior to border-adjacent --
this is not a difference-in-differences or switcher event study in the sense used elsewhere in this
research program. For farmland value, where we have three time points spanning a wide range of state
market ages (Colorado's market is 12 years old by 2026; Delaware's is under one year old), we
additionally estimate an event study by \emph{state market age} rather than calendar time: we interact
$\text{Border}_i$ with indicators for how many years have elapsed since county $i$'s own state began
retail sales, testing whether the border/interior gap changes as a state's legal market matures.

For crop substitution, Oregon's cultivation-license data (used in a companion cross-sectional check) has
no historical time dimension, but the border/interior treatment itself does not require it -- distance to
Idaho is a fixed geographic fact independent of any licensing dataset. We therefore estimate a
cross-sectional regression on Oregon's 36 counties,
\[
\Delta \ln(\text{cropland}_i) = \alpha + \beta \, \text{Border}_i + \gamma \ln(\text{population}_i) +
\delta \ln(\text{land area}_i) + \varepsilon_i,
\]
with $\Delta\ln(\text{cropland}_i)$ the change in log harvested cropland acres, 2012--2022.
\textbf{This is a single cross-section with one observation per county; there is no repeated
within-county observation to cluster on, so county-clustered standard errors are not applicable here.} We
report HC1 robust standard errors instead.

\textbf{Treatment and control, restated precisely:} in every specification in this paper, the
\emph{treatment} group is counties within 70 miles of a non-legal-state border, located inside a state
where recreational cannabis is already legal; the \emph{control} group is interior counties of the
\emph{same} legal states, more than 70 miles from any non-legal border. Neither group includes any county
in a non-legal state.

\section{Results}
\label{sec:results}

\subsection{Cannabis revenue: a clean, well-powered null}

Table~\ref{tab:revenue_border70} reports the revenue results. At the 70-mile threshold, Colorado
contributes 6 border counties and Washington 6, for 12 treated and 69 interior counties -- a substantial
improvement over a 50-mile threshold used in an earlier pass of this project, where Colorado's only
border county was Sedgwick alone, a small tourism-town outlier that drove a fragile, non-robust marginal
result. At 70 miles, the coefficient on Border is small, positive, and never close to significant
($b=+0.046$, county-clustered $p=0.83$ with full controls), and is essentially unchanged whether or not
Sedgwick is included. We read this as a genuine null, not an underpowered one: the treated group is
multi-county in composition in both states, and the point estimate itself is small.

\begin{table}[htbp]
\centering
\caption{Border-vs-Interior Estimates: Cannabis Retail Revenue per Capita (CO + WA)}
\label{tab:revenue_border70}
\begin{tabular}{lcccc}
\toprule
Specification & Coefficient & Std.\ Error (HC1) & SE (county-cl.) & $p$-value (county-cl.) \\
\midrule
No controls & $+0.051$ & $0.116$ & $0.260$ & $0.845$ \\
+ log(pop) + year FE + state FE & $+0.046$ & $0.090$ & $0.217$ & $0.833$ \\
\quad excluding Sedgwick County, CO & $-0.060$ & $0.083$ & $0.200$ & $0.763$ \\
\bottomrule
\end{tabular}
\begin{minipage}{0.95\textwidth}
\vspace{0.3em}
\footnotesize
\textit{Notes:} Outcome is log(cannabis retail revenue per capita). County-fiscal-year panel, Colorado
(2014--2025) and Washington (2015--2025), $n=815$. Treatment is a county within 70 miles of a non-legal
state; control is the interior counties of the same two states. At the 70-mile threshold Colorado
contributes 6 border counties (Larimer, Mesa, Montezuma, Morgan, Sedgwick, Weld) and Washington 6
(Asotin, Columbia, Pend Oreille, Spokane, Stevens, Whitman), for 12 border and 69 interior counties total
-- unlike the original 50-mile threshold used earlier in this project, where Colorado's only border
county was Sedgwick alone, a single-county outlier that drove an apparent (and non-robust) marginal
result. At 70 miles the treated group has genuine multi-county composition and the null is robust with or
without Sedgwick.
\end{minipage}
\end{table}


This null is not an average masking offsetting effects at different points in a state's market lifecycle:
Table~\ref{tab:revenue_percapita_eventstudy} and Figure~\ref{fig:revenue_percapita_eventstudy} extend it
into the same market-age event study used elsewhere in this paper, and no market-age bin approaches
significance. As with revenue per store, there are no pre-launch bins to plot -- no legal cannabis revenue
exists before a state's own retail market opens -- so this figure, like Figure~\ref{fig:revenue_perstore_eventstudy},
starts cleanly at market-age year 0.

\begin{table}[htbp]
\centering
\caption{Revenue-per-Capita Border Gap by State Market Age}
\label{tab:revenue_percapita_eventstudy}
\begin{tabular}{lcccc}
\toprule
Years since state's own retail launch & Border $\times$ Year coef. & SE (county-cl.) & $p$-value & Border obs.\\
\midrule
0 & $-0.421$ & $0.360$ & $0.243$ & 7 \\
1 & $+0.142$ & $0.226$ & $0.531$ & 7 \\
2 & $-0.287$ & $0.357$ & $0.422$ & 9 \\
3 & $-0.353$ & $0.343$ & $0.304$ & 11 \\
4 & $-0.168$ & $0.282$ & $0.551$ & 11 \\
5 & $-0.049$ & $0.279$ & $0.862$ & 11 \\
6 & $+0.085$ & $0.265$ & $0.750$ & 12 \\
7 & $+0.153$ & $0.288$ & $0.596$ & 12 \\
8 & $+0.304$ & $0.267$ & $0.255$ & 12 \\
9+ & $+0.334$ & $0.291$ & $0.251$ & 30 \\
\bottomrule
\end{tabular}
\begin{minipage}{0.95\textwidth}
\vspace{0.3em}
\footnotesize
\textit{Notes:} Outcome is log(revenue per capita); same specification and panel as
Table~\ref{tab:revenue_perstore_eventstudy}, $n=815$. No coefficient approaches significance at any
market age -- the per-capita null (Table~\ref{tab:revenue_border70}) holds throughout a state's market
lifecycle, not merely on average. Confidence intervals are wide and overlap zero in every bin, in sharp
contrast to Figure~\ref{fig:revenue_perstore_eventstudy}'s per-store premium, which is precisely the
contrast the license-cap mechanism in Section~\ref{sec:revenue_perstore} predicts: if the number of stores
cannot expand to meet elevated border demand, that demand should not show up when revenue is divided by
population, only when it is divided by the (capped) number of sellers.
\end{minipage}
\end{table}


\begin{figure}[htbp]
\centering
\includegraphics[width=0.95\textwidth]{figures/fig1_revenue_percapita_eventstudy.png}
\caption{Revenue per capita: border $\times$ market-age event-study coefficients, with 95\% confidence
intervals (county-clustered SEs). Compare with Figure~\ref{fig:revenue_perstore_eventstudy}: wide,
zero-crossing intervals throughout, versus a premium that becomes clearly distinguishable from zero by
market-age year 5.}
\label{fig:revenue_percapita_eventstudy}
\end{figure}

Since neither state's data extends before its own retail launch, there is no true pre-treatment period to
plot for either revenue outcome. The closest available substitute is to examine the control (interior)
group's own trajectory directly, stacking Colorado (launched 2014) and Washington (launched 2015) on a
common market-age clock rather than collapsing straight to a border-minus-interior difference.
Figure~\ref{fig:trends_percapita} plots border and interior counties' revenue-per-capita levels
separately: the two lines track each other closely across the entire market lifecycle, visually
confirming the null in Table~\ref{tab:revenue_border70} rather than merely reporting it as a coefficient
near zero.

\begin{figure}[htbp]
\centering
\includegraphics[width=0.95\textwidth]{figures/fig2_border_vs_interior_trends_percapita.png}
\caption{Revenue per capita: border and interior counties' trends plotted separately (not as a
difference), stacked on a common market-age clock. The two lines overlap throughout, the visual
counterpart of the null reported in Table~\ref{tab:revenue_border70} and Table~\ref{tab:revenue_percapita_eventstudy}.}
\label{fig:trends_percapita}
\end{figure}

\subsection{Revenue per store: the one significant result in this paper}
\label{sec:revenue_perstore}

Table~\ref{tab:revenue_perstore_border70} reports revenue per dispensary store, which tells a markedly
different story than revenue per capita. Border counties' stores generate significantly more revenue than
interior counties' stores: $b=+0.395$ (county-clustered $p=0.015$) with the full control set, corresponding
to roughly 48 percent higher revenue per store. This is not an artifact of the Sedgwick County outlier
discussed in Section~\ref{sec:data} -- excluding it, the estimate falls only modestly, to $+0.302$
($p=0.040$), and remains significant. Nor is it driven by one state: estimated separately, Colorado shows
$b=+0.552$ ($p=0.005$) and Washington $b=+0.441$ ($p=0.077$, marginal but directionally consistent),
despite the two states' raw (unconditional) county-level comparisons initially looking different --
Colorado's border stores average roughly double its interior stores' revenue, while Washington's border
and interior stores are closer in raw terms before controls.

This is, to our knowledge, the only genuinely significant result among the border-competition checks
conducted across this research program -- entry timing, entry density, and revenue per capita have all
come back null in various companion analyses. We think it resolves, rather than contradicts, those
earlier nulls. A companion dispensary-density analysis found no evidence that more stores open in border
counties, and proposed that state licensing caps on the number of dispensaries per county or population
could mask a real demand effect: if demand is elevated near a captive-market border but the number of
licensed stores cannot adjust to meet it, the excess demand should show up as more revenue per existing
store, not as more stores or higher aggregate per-capita revenue (which can stay flat if border counties
simply have fewer stores per capita, each doing more business). That is exactly the pattern in
Table~\ref{tab:revenue_perstore_border70}: a robust, two-state, outlier-independent revenue-per-store
premium alongside a null on both entry density and per-capita revenue.

\begin{table}[htbp]
\centering
\caption{Border-vs-Interior Estimates: Cannabis Revenue per Dispensary Store (CO + WA)}
\label{tab:revenue_perstore_border70}
\begin{tabular}{lcccc}
\toprule
Specification & Coefficient & SE (HC1) & SE (county-cl.) & $p$-value (county-cl.) \\
\midrule
No controls & $+0.502$ & $0.101$ & $0.233$ & $0.031$ \\
+ log(pop) + year FE + state FE & $+0.395$ & $0.074$ & $0.163$ & $0.015$ \\
\quad excluding Sedgwick County, CO & $+0.302$ & $0.068$ & $0.147$ & $0.040$ \\
\midrule
\multicolumn{5}{l}{\textit{Single-state estimates, + log(pop) + year FE (both HC1 $p<0.001$)}} \\
Colorado only (42 counties) & $+0.552$ & $<0.001$ & $0.199$ & $0.005$ \\
Washington only (29 counties) & $+0.441$ & $<0.001$ & $0.263$ & $0.077$ \\
\bottomrule
\end{tabular}
\begin{minipage}{0.95\textwidth}
\vspace{0.3em}
\footnotesize
\textit{Notes:} Outcome is log(total county cannabis revenue $\div$ number of currently-open dispensary
stores). Store counts are cumulative-open-by-year, built from geocoded dispensary license records
point-in-polygon assigned to counties, counted as open if the license's open date precedes each
fiscal/calendar year's own reporting cutoff (matching each state's revenue-reporting convention). Sample
restricted to county-years with at least one open store, $n=706$ (71 counties, 10 border and 61
interior). Unlike the revenue-per-capita and dispensary-density checks reported elsewhere in this
project -- both of which are null -- revenue per store is significantly \emph{higher} in border counties,
robust to dropping Sedgwick County (the single-county outlier that drove an earlier, non-robust
50-mile-threshold result) and to estimating each state separately. This is the only significant result
of the checks reported in this project's border-competition research thread, and it is directly
consistent with the license-cap explanation proposed when the original entry-density check came back null:
if border-county demand is genuinely elevated but the number of licensed stores is capped by state
regulation, the excess demand should appear as higher revenue per existing store rather than as more
stores or higher aggregate per-capita revenue -- exactly the pattern observed here.
\end{minipage}
\end{table}


\subsubsection{Event study by market age: the premium strengthens, it does not fade}

Table~\ref{tab:revenue_perstore_eventstudy} extends this result into an event study by state market age,
the same design used for farmland value in Section~\ref{sec:landvalue_marketage}. Market-age year 0 is too
thin to interpret (only 4 border observations, drawn from both states' earliest data), but from year 1
onward the coefficient is positive in \emph{every} bin -- no sign reversals -- and grows from marginal or
noisy in years 2--4 to significant and rising, from $+0.40$ at year 5 to $+0.59$ at year 9 and beyond. This
is the opposite pattern from the farmland-value event study, which turned negative only after several
years matured; here the premium strengthens rather than fades, consistent with a supply constraint
(license caps) that binds more tightly relative to demand as a legal market grows and matures, rather than
a one-time novelty effect that would be expected to fade.

\begin{table}[htbp]
\centering
\caption{Revenue-per-Store Border Premium by State Market Age}
\label{tab:revenue_perstore_eventstudy}
\begin{tabular}{lcccc}
\toprule
Years since state's own retail launch & Border $\times$ Year coef. & SE (county-cl.) & $p$-value & Border obs.\\
\midrule
0 & $-0.213$ & $0.563$ & $0.705$ & 4 \\
1 & $+0.296$ & $0.133$ & $0.026$ & 7 \\
2 & $+0.231$ & $0.164$ & $0.158$ & 8 \\
3 & $+0.131$ & $0.307$ & $0.670$ & 9 \\
4 & $+0.274$ & $0.188$ & $0.144$ & 9 \\
5 & $+0.397$ & $0.197$ & $0.044$ & 9 \\
6 & $+0.431$ & $0.185$ & $0.020$ & 10 \\
7 & $+0.563$ & $0.239$ & $0.019$ & 10 \\
8 & $+0.559$ & $0.234$ & $0.017$ & 10 \\
9+ & $+0.594$ & $0.230$ & $0.010$ & 26 \\
\bottomrule
\end{tabular}
\begin{minipage}{0.95\textwidth}
\vspace{0.3em}
\footnotesize
\textit{Notes:} Outcome is log(revenue per store); each row is the coefficient on
Border\_70mi $\times$ (market-age-year indicator), county-clustered SEs, controlling for log(population),
state fixed effects, and market-age-year fixed effects, $n=706$. Market-age year 0 is thin (only 4 border
observations, both states' earliest data) and should not be over-read. From year 1 onward the coefficient
is positive in every single bin -- no sign reversals -- and grows from marginal/noisy in years 2--4 to
significant and rising (+0.40 to +0.59) from year 5 onward. Unlike the farmland-value event study
(Table~\ref{tab:landvalue_marketage}), which turns negative only after several years, this pattern
strengthens rather than fades or reverses as a state's legal market matures, consistent with a license-cap
constraint that becomes more binding relative to demand over time.
\end{minipage}
\end{table}


\begin{figure}[htbp]
\centering
\includegraphics[width=0.95\textwidth]{figures/fig3_revenue_perstore_eventstudy.png}
\caption{Revenue per store: border $\times$ market-age event-study coefficients, with 95\% confidence
intervals (county-clustered SEs). Unlike Figure~\ref{fig:landvalue_eventstudy}, there is no pre-launch
period to plot: no legal cannabis revenue exists before a state's own retail market opens, so this
coefficient path starts cleanly at year 0.}
\label{fig:revenue_perstore_eventstudy}
\end{figure}

As with revenue per capita, Figure~\ref{fig:trends_perstore} plots border and interior counties' levels
separately rather than only their difference, using the same stacked market-age clock. The two groups
start close together in a market's first two years -- both near their lowest levels, since the whole
market is still small -- and then visibly diverge from year 2 onward, with border counties pulling ahead
and staying ahead through year 9+. Because neither state has data before its own launch, we cannot show
the two lines were on parallel trajectories \emph{before} treatment in the way Figure~\ref{fig:landvalue_eventstudy}
does for farmland value; what this figure does show is that the two groups are not visibly different from
the earliest available data and diverge only afterward, which is the most this specific data can offer in
place of a true pre-trend check.

\begin{figure}[htbp]
\centering
\includegraphics[width=0.95\textwidth]{figures/fig4_border_vs_interior_trends_perstore.png}
\caption{Revenue per store: border and interior counties' trends plotted separately (not as a
difference), stacked on a common market-age clock. The two groups start together and diverge from year 2
onward -- the visual counterpart of the growing premium in Table~\ref{tab:revenue_perstore_eventstudy}.}
\label{fig:trends_perstore}
\end{figure}

\subsubsection{Mechanism checks: does the premium behave like genuine cross-border demand?}
\label{sec:mechanism}

The results above are consistent with the captive cross-border demand hypothesis but do not directly test
it: an observed border/interior gap in revenue per store could in principle reflect some other difference
between border and interior counties rather than cross-border shoppers specifically. We probe the
mechanism two further ways.

First, a \emph{dose-response} check: if the effect reflects shoppers traveling from the neighboring
non-legal state, revenue per store should rise smoothly as a county gets closer to that state's border, not
merely jump at the 70-mile threshold used elsewhere in this paper. We replace the binary indicator with
log(distance to the nearest non-legal-state county) and re-estimate the same specification. The
coefficient has the expected sign (closer to the border, more revenue per store: $b=-0.182$) and is
significant under heteroskedasticity-robust standard errors ($p=0.001$), but loses significance once
standard errors are clustered by county ($p=0.124$), and weakens further excluding Sedgwick County
($p=0.351$ clustered). This is suggestive of a genuine distance gradient but not, on its own, strong
confirmatory evidence. Figure~\ref{fig:doseresponse_distance} plots this relationship directly as a
binned scatter of revenue per store against distance, both residualized on the same controls
(population, year, and state fixed effects) as the regression: the closest bin sits well above the fitted
line and the two farthest bins sit well below it, consistent with the estimated slope, but the middle of
the distribution is noisy and not monotonic, visually matching why the relationship does not survive
clustering.

\begin{figure}[htbp]
\centering
\includegraphics[width=0.85\textwidth]{figures/fig5_doseresponse_distance.png}
\caption{Dose-response: revenue per store against log(distance to non-legal border), both residualized
on log(population), year fixed effects, and state fixed effects. Points are 10-bin (decile) means of the
residualized distance variable; point size reflects the number of observations in each bin. The fitted
line's slope matches the regression coefficient in Table~\ref{tab:revenue_perstore_mechanism} exactly, by
construction.}
\label{fig:doseresponse_distance}
\end{figure}

Second, a \emph{captive market size} check, which we view as a better-motivated test of the mechanism than
distance to the nearest non-legal county alone: a county next to one small non-legal county and a county
next to a large non-legal metropolitan area are treated identically by a nearest-county distance measure,
but only the latter has a large population of potential cross-border shoppers to draw from. We construct,
for each Colorado and Washington county, the total population of \emph{all} non-legal-state counties within
70 miles -- not merely the nearest one -- via a radius query against the full national county population
distribution, and replace the distance measure with log(captive non-legal population). This measure is
significantly positive even under county-clustered standard errors ($b=+0.035$, $p=0.014$), and remains
significant excluding Sedgwick ($b=+0.027$, $p=0.034$). Only 10 of the 71 Colorado and Washington counties
in this panel have any non-legal population within 70 miles at all, so this check is still identified from
a similar small set of border counties as the rest of this section -- it is not new geographic variation --
but it reweights that same set by the actual size of the market each county could plausibly absorb, and it
is the most robust mechanism evidence in this paper.

\begin{table}[htbp]
\centering
\caption{Mechanism Checks for the Revenue-per-Store Border Premium}
\label{tab:revenue_perstore_mechanism}
\begin{tabular}{lcccc}
\toprule
Specification & Coefficient & SE (HC1) & SE (county-cl.) & $p$-value (county-cl.) \\
\midrule
\multicolumn{5}{l}{\textit{(a) Dose-response: log(distance to non-legal border), full controls}} \\
Full sample & $-0.182$ & $0.054$ & $0.118$ & $0.124$ \\
Excluding Sedgwick County, CO & $-0.107$ & $0.052$ & $0.114$ & $0.351$ \\
\midrule
\multicolumn{5}{l}{\textit{(b) Captive market size: log(non-legal population within 70mi), full controls}} \\
Full sample & $+0.035$ & $0.007$ & $0.014$ & $0.014$ \\
Excluding Sedgwick County, CO & $+0.027$ & $0.006$ & $0.013$ & $0.034$ \\
\bottomrule
\end{tabular}
\begin{minipage}{0.95\textwidth}
\vspace{0.3em}
\footnotesize
\textit{Notes:} Outcome is log(revenue per store), same CO+WA panel as Table~\ref{tab:revenue_perstore_border70},
$n=706$. Panel (a) replaces the binary border indicator with log(distance in miles to the nearest
non-legal-state county); a negative coefficient indicates revenue per store rises as a county gets closer
to a non-legal border, a smoother version of the binary result. Panel (b) replaces distance-to-nearest-county
with a market-size measure: the total population of \emph{all} non-legal-state counties within 70 miles
(not merely the nearest one), computed via a radius query against the full national county population
distribution. Only 10 of 71 CO/WA counties have any non-legal population within 70 miles at all, so both
checks are identified from a similar small set of border counties, now weighted by the actual size of the
captive market rather than treated as a uniform category. The dose-response check has the theoretically
expected sign but loses significance once standard errors are clustered by county; the captive-population
check is significant even under clustering, with or without Sedgwick, and is the most robust support for
the demand-absorption mechanism found in this paper.
\end{minipage}
\end{table}


\begin{figure}[htbp]
\centering
\includegraphics[width=0.85\textwidth]{figures/fig6_doseresponse_captivepop.png}
\caption{Dose-response: revenue per store against log(captive non-legal population within 70 miles), both
residualized on log(population), year fixed effects, and state fixed effects. The x-axis makes the
identifying variation visible directly: the large majority of counties cluster near zero captive
population (61 of 71 have none at all), and the apparent upward slope is carried by only two
distinctly-positioned bins at the right -- this is a real, significant relationship, but it rests on a
handful of high-captive-population counties, not a smooth gradient sampled across many distinct exposure
levels.}
\label{fig:doseresponse_captivepop}
\end{figure}

Taken together, the dose-response and captive-population checks point the same direction as the main
result but differ in how much weight they can bear on their own: the simple distance gradient is fragile
to clustering, while the population-weighted market-size measure survives it. We read this as meaningful
support for -- not proof of -- the cross-border demand interpretation, and note explicitly what would make
the case stronger: direct data on customer origin (out-of-state identification records, if a state
regulator were to publish them) would test the mechanism far more directly than any distance- or
population-based proxy can.

\subsubsection{Placebo check: the premium does not show up in general retail}
\label{sec:placebo}

A different concern than mechanism is specification: even a robust, well-identified border/interior gap in
cannabis revenue per store could reflect something about border counties in general -- through-traffic,
tourism, price arbitrage on goods generally -- rather than anything specific to cannabis and its captive
non-legal-market demand. We test this directly using Colorado's general retail sales data
\citep{CDORRetailByCounty} as a placebo outcome, applying the identical border/interior specification to
all-industry retail sales per capita instead of cannabis revenue. Table~\ref{tab:placebo_general_retail}
reports the result: a precise null ($b=-0.006$, $p=0.95$, county-clustered, $n=110$, 2016--2025). General
retail activity in Colorado's border counties is statistically indistinguishable from its interior
counties', even as cannabis revenue per store is significantly higher in the same border counties over the
same period. This is, in our view, the single most important robustness check in this paper: it directly
rules out the most obvious alternative explanation a referee would raise, that the revenue-per-store
premium is a generic border-town phenomenon rather than a cannabis-specific captive-demand effect. The
caveat, noted in Section~\ref{sec:data}, is that Colorado's disclosure rules limit this placebo's border
group to 3 of the main analysis's 6 border counties (Larimer, Mesa, and Weld), so it is not a perfect
match -- but it is the best available test, and it points the right direction.

\begin{table}[htbp]
\centering
\caption{Placebo Check: General Retail Sales per Capita, Colorado}
\label{tab:placebo_general_retail}
\begin{tabular}{lcccc}
\toprule
Specification & Coefficient & SE (HC1) & SE (county-cl.) & $p$-value (county-cl.) \\
\midrule
No controls & $-0.043$ & $0.053$ & $0.069$ & $0.533$ \\
+ log(pop) + year FE & $-0.006$ & $0.028$ & $0.087$ & $0.947$ \\
\bottomrule
\end{tabular}
\begin{minipage}{0.95\textwidth}
\vspace{0.3em}
\footnotesize
\textit{Notes:} Outcome is log(general retail sales per capita, all industries), Colorado Department of
Revenue, county-by-industry monthly filings aggregated to complete calendar years, 2016--2025,
$n=110$. CDOR discloses this industry-level breakdown only for Colorado's 11 largest counties by
filing volume (Adams, Arapahoe, Boulder, Denver, Douglas, El Paso, Jefferson, Larimer, Mesa, Pueblo,
Weld), which includes 3 of the 6 border counties in the main cannabis-revenue analysis (Larimer, Mesa,
Weld) and none of the smaller ones (Montezuma, Morgan, Sedgwick). County-clustered SEs. The border
coefficient is a precise, well-powered null -- general retail activity does not respond to proximity to a
non-legal state the way cannabis revenue per store does, supporting a cannabis-specific rather than a
generic border-town-shopping interpretation of the main result.
\end{minipage}
\end{table}


\subsubsection{Suggestive evidence on the mechanism from Washington's 2014 license lottery}
\label{sec:lottery}

One further piece of evidence, offered as suggestive context rather than a formal empirical test, comes
from Washington's initial 2014 licensing round. Unlike Colorado, Washington capped the statewide number of
retail licenses at 334 from the outset and allocated them by lottery in any jurisdiction where applications
exceeded the local quota \citep{WSLCB2014LotteryPR}. Statewide, more than 2,000 applicants competed for
those 334 licenses. Contemporary local reporting on the allocation gives a directly relevant data point:
Spokane County -- home to Spokane, Spokane Valley, and rural Spokane County, three of the jurisdictions
that make up one of this paper's six border counties -- received 143 applications against an 18-license
county allocation, with the City of Spokane itself drawing 85 applications for 8 available store licenses,
Spokane Valley 41 applications for 3, and the remainder of the county 17 applications for 7
\citep{SpokesmanReview2014Lottery}. We do not build a formal test from this: a complete, comparable
applicant-to-license ratio for every border and interior jurisdiction in Washington was never compiled into
a single official dataset, and reconstructing one from scattered contemporary news coverage across the
whole state was beyond the scope of this paper. We report the Spokane figures because they are a
concrete, well-documented instance of exactly the mechanism this paper argues for -- real, oversubscribed
demand for cannabis retail entry in a border county, resolved by a binding cap on the number of licenses
rather than by the market clearing through additional entry -- consistent with, though not proof of, the
revenue-per-store premium's interpretation in Section~\ref{sec:revenue_perstore}.

\subsection{Farmland value: null in levels, but with two results worth flagging}
\label{sec:landvalue_marketage}

Table~\ref{tab:landvalue_border_panel} reports the farmland-value panel. The raw and lightly controlled
specifications show a large, highly significant positive border/interior gap ($b=+0.31$, $p<0.001$), but
this is entirely a between-state artifact: some states simply have systematically higher farmland values
everywhere, for reasons unrelated to any individual county's border proximity. Once state fixed effects
are added, the gap collapses to a precise, insignificant near-zero ($b=-0.005$, $p=0.87$), and stays there
in the subsample restricted to the 17 states with genuine border-county variation.

Two further results complicate a simple ``null'' reading. First, the continuous distance version --
replacing the binary threshold with log-distance to the nearest non-legal border -- remains significant
with the full control set, but in the \emph{opposite} direction: greater distance from a non-legal border
is associated with \emph{higher}, not lower, land value ($b=+0.083$, $p=0.001$). This is the opposite sign
from what the captive-demand hypothesis predicts, and it disagrees with the binary specification's null.
We do not resolve this discrepancy in favor of either specification; we report both, since specification
sensitivity of this kind is itself informative about how fragile any land-value/border relationship is in
this data.

\begin{table}[htbp]
\centering
\caption{Border-vs-Interior Panel Estimates: Farmland Value per Acre, 21 Legal States}
\label{tab:landvalue_border_panel}
\begin{tabular}{lcccc}
\toprule
Specification & Coefficient & SE (HC1) & SE (county-cl.) & $p$-value (county-cl.) \\
\midrule
No controls & $+0.310$ & $0.036$ & $0.059$ & $<0.001$ \\
+ log(pop) & $+0.335$ & $0.030$ & $0.048$ & $<0.001$ \\
+ year FE & $+0.310$ & $0.035$ & $0.059$ & $<0.001$ \\
+ state FE & $-0.035$ & $0.028$ & $0.044$ & $0.428$ \\
+ log(pop) + year FE + state FE & $+0.008$ & $0.024$ & $0.039$ & $0.836$ \\
+ log(pop) + land area + year FE + state FE & $-0.005$ & $0.020$ & $0.032$ & $0.869$ \\
\midrule
\multicolumn{5}{l}{\textit{Continuous distance version (full controls), log(miles to non-legal border)}} \\
+ log(pop) + land area + year FE + state FE & $+0.083$ & $0.015$ & $0.025$ & $0.001$ \\
\midrule
\multicolumn{5}{l}{\textit{17-state subsample, dropping AZ/CA/NV/RI (no county within 70mi of a non-legal border)}} \\
+ log(pop) + land area + year FE + state FE & $-0.006$ & $0.020$ & $0.032$ & $0.847$ \\
\bottomrule
\end{tabular}
\begin{minipage}{0.95\textwidth}
\vspace{0.3em}
\footnotesize
\textit{Notes:} Outcome is log(farmland asset value per acre), NASS Census of Agriculture, 2012/2017/2022.
County-year panel, 947 counties matched across 21 states with an operating recreational retail market as
of 2026-08-05 and available land-value data, $n=2{,}835$. Treatment (border\_70mi) is 1 for counties
within 70 miles of a non-legal-state border, 0 otherwise (434 border counties, 513 interior). Standard
errors clustered by county. The raw and lightly controlled specifications show a large, significant
positive gap that is entirely a between-state artifact: it collapses to a precise null once state fixed
effects are added, and stays null in the subsample restricted to states with genuine border-county
variation. The continuous log-distance version, however, remains significant with the full control set --
in the \emph{opposite} direction of the binary specification (farther from a non-legal border associated
with \emph{higher}, not lower, land value) -- a specification-sensitivity result discussed in Section
\ref{sec:results}, not resolved in favor of either specification.
\end{minipage}
\end{table}


Second, Table~\ref{tab:landvalue_marketage} and Figure~\ref{fig:landvalue_eventstudy} report an event
study by state market age, including years \emph{before} a state's own legalization, which we can observe
because NASS's 2012/2017/2022 Census waves predate several states' launch dates entirely -- all three
waves predate Delaware's 2025 launch, and two of the three predate Illinois's 2020 launch. An earlier
version of this check discarded these pre-launch observations under the mistaken assumption that none
existed; once included, they reveal a complication that changes how the post-launch pattern should be
read. The border/interior gap is not flat before legalization: it is significantly \emph{negative} at
2 years before launch ($b=-0.111$, $p=0.031$) and significantly \emph{positive} at 1 year before launch
($b=+0.119$, $p=0.009$), fluctuations comparable in size to anything observed after legalization. The
immediate transition across the launch date itself is unremarkable -- the year-before and year-of-launch
coefficients are nearly identical ($+0.119$ and $+0.114$) -- but the pre-launch period as a whole is not the
flat, quiet baseline a clean causal design would want. Only starting around 2 years after launch does the
coefficient turn consistently negative and grow in magnitude, reaching $-0.384$ ($p=0.008$) at 7 years and
$-0.390$ ($p=0.002$) at 8-plus years.

We do not think this pattern can be read as a confident causal effect of legal-market access on farmland
values. The significant pre-launch swings mean we cannot rule out that border counties in states that
later legalize were already on a noisier or systematically different land-value trajectory than interior
counties, for reasons unrelated to cannabis policy -- exactly the kind of parallel-trends violation this
research program has flagged transparently elsewhere rather than smoothing over. What we can say is that
the specific decline visible from 2 years post-launch onward does not have an obvious pre-launch analogue
of the same sustained, growing, one-directional shape; the pre-launch swings are volatile and change sign,
while the post-launch decline is more gradual and monotonic. We report both facts and let the reader weigh
them, rather than asserting the post-launch pattern is definitively either a real effect or pure
continuation of pre-existing noise. Its direction is nonetheless worth noting because it runs opposite to
the revenue-per-store premium's own market-age pattern (Section~\ref{sec:revenue_perstore}), which strengthens
rather than fades or reverses over the same kind of horizon and has no pre-launch period to complicate it
in the same way, since no legal cannabis revenue exists before a state's own market opens.

\begin{table}[htbp]
\centering
\caption{Land-Value Border Gap by State Market Age, Including Pre-Launch Years}
\label{tab:landvalue_marketage}
\begin{tabular}{lcccc}
\toprule
Years relative to state's own retail launch & Border $\times$ Year coef. & SE (county-cl.) & $p$-value & $n$ \\
\midrule
\multicolumn{5}{l}{\textit{Pre-launch (no legal cannabis market yet exists in the state)}} \\
$-4$ or earlier & $+0.054$ & $0.036$ & $0.129$ & 1{,}324 \\
$-3$ & $-0.091$ & $0.050$ & $0.069$ & 157 \\
$-2$ & $-0.111$ & $0.051$ & $0.031$ & 273 \\
$-1$ & $+0.119$ & $0.046$ & $0.009$ & 296 \\
\midrule
\multicolumn{5}{l}{\textit{Post-launch}} \\
0 & $+0.114$ & $0.101$ & $0.263$ & 204 \\
2 & $-0.119$ & $0.057$ & $0.035$ & 154 \\
3 & $-0.144$ & $0.073$ & $0.051$ & 185 \\
4 & $-0.315$ & $0.158$ & $0.046$ & 72 \\
7 & $-0.384$ & $0.144$ & $0.008$ & 36 \\
8+ & $-0.390$ & $0.129$ & $0.002$ & 102 \\
\bottomrule
\end{tabular}
\begin{minipage}{0.95\textwidth}
\vspace{0.3em}
\footnotesize
\textit{Notes:} Outcome is log(farmland value per acre); each row is the coefficient on
Border\_70mi $\times$ (market-age-year indicator) from a single regression with county-clustered SEs,
controlling for log(population), log(land area), state fixed effects, and market-age-year fixed effects.
Market-age year 1 and year 5 are omitted -- state fixed effects fully absorb the border indicator in those
bins (only one state contributes observations at that specific age), leaving the interaction
unidentified rather than a genuine zero. Pre-launch years use NASS Census-of-Agriculture waves (2012,
2017, 2022) that predate a state's own legalization -- e.g. all three waves predate Delaware's 2025
launch, and 2012/2017 predate Illinois's 2020 launch -- so these are not a placebo constructed after the
fact; the underlying land-value data genuinely exists before treatment for many states in this panel. See
Figure~\ref{fig:landvalue_eventstudy} for the full coefficient path and Section~\ref{sec:landvalue_marketage}
for why the pre-launch pattern complicates a simple causal reading of the post-launch decline.
\end{minipage}
\end{table}


\begin{figure}[htbp]
\centering
\includegraphics[width=0.95\textwidth]{figures/fig7_landvalue_eventstudy.png}
\caption{Land value: border $\times$ market-age event-study coefficients, with 95\% confidence intervals
(county-clustered SEs). The dashed vertical line marks each state's own retail launch. Years 1 and 5 are
omitted (unidentified -- state fixed effects fully absorb the border indicator when only one state
contributes observations at that age).}
\label{fig:landvalue_eventstudy}
\end{figure}

Figure~\ref{fig:trends_landvalue} shows the same information a different way: border and interior
counties' levels plotted separately, including the genuine pre-launch years, rather than collapsed into a
single difference. Both groups trend upward together through the pre-launch period, with border running
consistently a little above interior throughout -- a small, fairly stable gap that predates legalization
entirely and is visible directly in the figure, not just inferable from the significant pre-launch
coefficients in Table~\ref{tab:landvalue_marketage}. After launch, interior counties continue climbing
(reaching their highest levels in the panel by 7--8 years out) while border counties' values flatten and
stop growing. The post-launch decline reported above is thus better described as interior counties
pulling further ahead of an already-existing gap, not border counties falling in absolute terms -- a more
precise, and more clearly non-causal-looking, characterization than the coefficient table alone conveys.

\begin{figure}[htbp]
\centering
\includegraphics[width=0.95\textwidth]{figures/fig8_border_vs_interior_trends_landvalue.png}
\caption{Land value: border and interior counties' trends plotted separately (not as a difference),
including real pre-launch years. The dashed vertical line marks each state's own retail launch. Both
groups trend upward together before launch, with border running consistently a little above interior;
after launch, interior continues climbing while border flattens.}
\label{fig:trends_landvalue}
\end{figure}

\subsection{Crop substitution: underpowered to the point of being uninformative}

Table~\ref{tab:cropsub_border70} reports the Oregon cross-section. Oregon's only non-legal-state neighbor
is Idaho, and only 3 of its 36 counties -- Baker, Malheur, and Wallowa -- fall within 70 miles of that
border. With a treated group of three counties, no specification approaches significance, and the point
estimates are not even consistently signed across the three counties themselves (Baker's cropland grew
3.7 percent over 2012--2022; Malheur's fell 21.7 percent; Wallowa's fell 10.3 percent). We do not read
this as evidence against a border effect on crop substitution; the design simply lacks the statistical
power to detect one in Oregon. We note, in addition, that Jackson and Josephine counties -- which have by
far the largest cannabis cultivation footprint of any Oregon county in a companion license-based analysis
-- are 313 and 357 miles from Idaho, situated instead near the California border (itself a legal state).
Oregon's actual cannabis cultivation activity is concentrated nowhere near its captive-demand border; it
sits in the historic Emerald-Triangle-adjacent production region in the state's southwest, a legacy of
decades of illicit cultivation predating legalization, not a response to cross-border demand. This is a
useful negative fact in its own right: in Oregon at least, cultivation intensity and non-legal-border
proximity are not the same geography, and a researcher using one as a proxy for the other would be
measuring something different than intended.

\begin{table}[htbp]
\centering
\caption{Border-vs-Interior Estimates: Harvested Cropland Change, Oregon}
\label{tab:cropsub_border70}
\begin{tabular}{lcccc}
\toprule
Specification & Coefficient & SE (HC1) & $p$-value & $R^2$ \\
\midrule
border\_70mi, no controls & $-0.045$ & $0.078$ & $0.560$ & $0.004$ \\
border\_70mi + log(pop) + log(land area) & $-0.038$ & $0.100$ & $0.706$ & $0.007$ \\
log(distance to non-legal border), no controls & $+0.026$ & $0.049$ & $0.587$ & $0.006$ \\
log(distance) + log(pop) + log(land area) & $+0.022$ & $0.066$ & $0.738$ & $0.007$ \\
\bottomrule
\end{tabular}
\begin{minipage}{0.95\textwidth}
\vspace{0.3em}
\footnotesize
\textit{Notes:} Outcome is log change in NASS harvested cropland acres, 2012--2022, all 36 Oregon
counties. Cross-sectional, HC1 robust SEs (county clustering not applicable -- one observation per
county). Oregon's only non-legal-state neighbor is Idaho; only 3 of 36 counties (Baker, Malheur, Wallowa)
fall within 70 miles of it, so the treated group is extremely thin and no specification approaches
significance. Notably, Jackson and Josephine counties -- which had by far the largest cannabis
cultivation canopy in a companion license-based analysis -- are 313 and 357 miles from Idaho, near the
California border instead (itself a legal state): the counties with the most cannabis cultivation activity
in Oregon are not the border counties this design treats, underscoring that captive-demand-border
proximity and cultivation intensity are empirically distinct phenomena in this state.
\end{minipage}
\end{table}


\subsubsection{An event-study view, for comparability -- and why it should not be over-read}

The other three outcomes in this paper are shown as event studies because either multiple states could be
stacked on a common market-age clock (revenue) or a single state's own panel spans years before and after
its launch (farmland value). Oregon offers neither advantage for crop substitution: it is the only state
with usable cultivation-border data, so there is no second state to stack against it, though its 3 NASS
waves (2012, 2017, 2022) do span Oregon's own 2015 launch, giving three points at event years $-3$, $+2$,
and $+7$. Table~\ref{tab:cropsub_eventstudy} and Figure~\ref{fig:cropsub_eventstudy} report it for
comparability with the rest of the paper, but it should not be read the same way. The coefficient here is
estimated on cropland \emph{levels}, not the year-over-year change used in the main cross-section, and it
comes back nearly identical at all three event years -- $+0.78$ in 2012, $+0.81$ in 2017, $+0.73$ in 2022 --
\emph{including three years before Oregon had any legal cannabis market at all}. Figure~\ref{fig:cropsub_trends}
shows why directly: border and interior counties' cropland levels run on two flat, parallel, non-converging
lines throughout the entire period. This is not a treatment effect; it is a fixed characteristic of which
counties happen to be near Idaho -- Baker, Malheur, and Wallowa are large wheat-growing counties in eastern
Oregon, larger in cropland terms than most interior counties regardless of cannabis policy. It is precisely
why the main analysis in Table~\ref{tab:cropsub_border70} uses the \emph{change} in cropland rather than its
level: differencing removes this stable gap, which is exactly what makes the level comparison uninformative
as a treatment estimate even though it looks superficially like a large, marginally significant effect.

\begin{table}[htbp]
\centering
\caption{Crop Substitution Event Study, Oregon Only (Illustrative)}
\label{tab:cropsub_eventstudy}
\begin{tabular}{lcccc}
\toprule
Event year (calendar year) & Border $\times$ Year coef. & SE (county-cl.) & $p$-value \\
\midrule
$-3$ (2012) & $+0.776$ & $0.425$ & $0.068$ \\
$+2$ (2017) & $+0.812$ & $0.430$ & $0.059$ \\
$+7$ (2022) & $+0.731$ & $0.393$ & $0.063$ \\
\bottomrule
\end{tabular}
\begin{minipage}{0.95\textwidth}
\vspace{0.3em}
\footnotesize
\textit{Notes:} Outcome is log(harvested cropland acres), \emph{levels}, not the change used in
Table~\ref{tab:cropsub_border70}. $n=108$ (36 Oregon counties $\times$ 3 NASS years), only 3 border
counties (Baker, Malheur, Wallowa) at every year, so treat as illustrative rather than well-powered.
Unlike the other three outcomes in this paper, Oregon is the only state with usable crop-substitution
data, so there is no cross-state market-age clock to stack on -- this event study uses Oregon's own 3
NASS waves (2012, 2017, 2022) relative to its 2015 retail launch. The coefficient is nearly identical
across all three years, \emph{including 2012, three years before Oregon legalized}: border counties simply
have more harvested cropland than interior counties in level terms throughout the entire period (Baker,
Malheur, and Wallowa are large wheat-growing counties in eastern Oregon), a fixed characteristic unrelated
to legalization timing. This is precisely why the main cross-sectional analysis
(Table~\ref{tab:cropsub_border70}) uses the \emph{change} in cropland rather than its level: differencing
removes this stable level gap, which is exactly what makes it not the right measure for a treatment effect.
\end{minipage}
\end{table}


\begin{figure}[htbp]
\centering
\includegraphics[width=0.8\textwidth]{figures/fig9_cropsub_eventstudy.png}
\caption{Crop substitution event study, Oregon only (illustrative -- only 3 border counties at every
year). Outcome is log(cropland) in \emph{levels}. The coefficient is essentially flat across all three
years, including 2012, three years before Oregon's own legalization.}
\label{fig:cropsub_eventstudy}
\end{figure}

\begin{figure}[htbp]
\centering
\includegraphics[width=0.85\textwidth]{figures/fig10_cropsub_trends.png}
\caption{Crop substitution: border and interior counties' cropland levels plotted separately, Oregon
only. The two lines run flat and parallel throughout, including the pre-launch year -- visual confirmation
that the gap in Table~\ref{tab:cropsub_eventstudy} is a fixed characteristic of these particular counties,
not a legalization-related effect.}
\label{fig:cropsub_trends}
\end{figure}

\subsection{Extending the design: Oregon and Montana}
\label{sec:or_mt}

Colorado and Washington's revenue-per-store result raises an obvious question: does it replicate in other
states with genuinely independent data sources? Oregon and Montana both border non-legal states (Idaho
for Oregon; Idaho, North Dakota, South Dakota, and Wyoming for Montana -- Montana is the only state in
this design bordering more than one), and, as Section~\ref{sec:data_revenue} describes, both turned out to
have usable regulator revenue data once we looked past the states with statewide-only or dashboard-only
reporting. We treat this as an out-of-sample extension of the main result, not a re-specification: the
treatment/control definition, control set, and estimating equation are unchanged.

\begin{table}[htbp]
\centering
\caption{Border-vs-Interior Estimates: Revenue per Store, Pooled Colorado + Washington + Oregon}
\label{tab:pooled_perstore_or}
\begin{tabular}{lcccc}
\toprule
Specification & Coefficient & SE (HC1) & SE (county-cl.) & $p$-value (county-cl.) \\
\midrule
No controls & $+0.526$ & $0.099$ & $0.226$ & $0.020$ \\
+ log(pop) + year FE + state FE & $+0.429$ & $0.074$ & $0.157$ & $0.006$ \\
\quad excluding Sedgwick, CO and Malheur, OR & $+0.323$ & $0.067$ & $0.142$ & $0.023$ \\
\midrule
\multicolumn{5}{l}{\textit{Single-state estimates, + log(pop) [+ year FE for CO/WA] (all HC1 $p<0.01$)}} \\
Colorado only & $+0.552$ & $0.088$ & -- & -- \\
Washington only & $+0.441$ & $0.101$ & -- & -- \\
Oregon only (2025 cross-section) & $+1.168$ & $0.433$ & -- & -- \\
\bottomrule
\end{tabular}
\begin{minipage}{0.95\textwidth}
\vspace{0.3em}
\footnotesize
\textit{Notes:} Adds Oregon to the Colorado+Washington revenue-per-store panel in
Table~\ref{tab:revenue_perstore_border70}. Oregon's county sales come from OLCC's own Socrata series
(2017--2025, all years complete); Oregon's store counts, unlike Colorado's and Washington's, cannot be
built as a historical cumulative-open panel -- the geocoded dispensary listing used elsewhere in this
project has real coverage gaps for Oregon specifically (several records geocode outside Oregon entirely,
and Baker County, one of the three border counties, has zero matches despite recording revenue in every
year 2017--2025) -- so Oregon instead contributes a single 2025 cross-section, pairing 2025 revenue with
currently-active retailer counts from OLCC's own license registry (which tags county natively). Pooling
this cross-section into the panel is valid despite the year mismatch: year fixed effects absorb any
Oregon-specific 2025 level shock the same way they absorb Colorado's and Washington's own year effects,
and state fixed effects absorb any state-level difference in revenue per store, so Oregon's within-state
border/interior contrast is identified the same way Colorado's and Washington's are. Adding Oregon
\emph{strengthens} the pooled estimate's significance under county-clustered SEs ($p=0.015\to p=0.006$),
and all three states are independently positive and significant. Malheur County, OR (Ontario, directly
across the Snake River from Fruitland/Payette, Idaho) plays the same outlier role for Oregon that Sedgwick
County plays for Colorado; excluding both, the pooled estimate remains significant.
\end{minipage}
\end{table}


Pooling Oregon into the Colorado-Washington revenue-per-store panel (Table~\ref{tab:pooled_perstore_or})
\emph{strengthens} the result: the county-clustered $p$-value falls from $0.015$ to $0.006$, and all three
states are independently positive and significant, despite Oregon's revenue-per-store outcome coming from
a different underlying construction (a single-year cross-section against currently active licenses, rather
than a historical panel of cumulative store openings) than Colorado's and Washington's. Oregon has its own
version of the Sedgwick County pattern: Malheur County -- home to Ontario, Oregon, a border town directly
across the Snake River from Fruitland and Payette, Idaho, and a well-known cross-border cannabis
destination in its own right -- generates far more revenue per store than any interior Oregon county, but
excluding it leaves the pooled estimate significant, just as excluding Sedgwick does for Colorado and
Washington alone.

\begin{table}[htbp]
\centering
\caption{Oregon and Montana Standalone: Revenue per Capita and Revenue per Store}
\label{tab:or_mt_standalone}
\begin{tabular}{lcccc}
\toprule
Specification & Coefficient & SE (HC1) & SE (county-cl.) & $p$-value (county-cl.) \\
\midrule
\multicolumn{5}{l}{\textit{Oregon -- outcome: log(revenue per capita), 2017--2025 panel, $n=255$}} \\
+ log(pop) + year FE & $+0.968$ & $0.317$ & $0.667$ & $0.147$ \\
\quad excluding Malheur County & $+0.233$ & $0.279$ & $0.401$ & $0.562$ \\
\midrule
\multicolumn{5}{l}{\textit{Oregon -- outcome: log(revenue per store), 2025 cross-section, $n=30$ (HC1 only)}} \\
+ log(pop) & $+1.168$ & $0.433$ & -- & $p=0.007$ \\
\quad excluding Malheur County & $+0.669$ & $0.194$ & -- & $p=0.001$ \\
\midrule
\multicolumn{5}{l}{\textit{Montana -- outcome: log(revenue per capita), 2022--2025 panel, $n=117$}} \\
+ log(pop) + year FE & $+0.654$ & $0.180$ & $0.311$ & $0.036$ \\
\quad excluding Richland County & $+0.531$ & $0.177$ & $0.304$ & $0.081$ \\
\midrule
\multicolumn{5}{l}{\textit{Montana -- outcome: log(revenue per store), 2025 cross-section, $n=23$ (HC1 only)}} \\
+ log(pop) & $+0.012$ & $0.232$ & -- & $p=0.958$ \\
\bottomrule
\end{tabular}
\begin{minipage}{0.95\textwidth}
\vspace{0.3em}
\footnotesize
\textit{Notes:} Oregon: 3 border counties (Baker, Malheur, Wallowa, all $\leq$70mi from Idaho); revenue
from OLCC's Socrata series, store counts from OLCC's license registry (2025 snapshot; see
Table~\ref{tab:pooled_perstore_or}'s notes). Montana: 11 border counties with nonzero recorded cannabis
revenue ($\leq$70mi from Idaho, North Dakota, South Dakota, or Wyoming -- Montana is the only state in
this design bordering more than one non-legal state); revenue scraped from the Department of Revenue's
monthly county sales reports (revenue.mt.gov, no bulk download exists, 53 of 54 months from January 2022
through June 2026 retrieved directly, one retried successfully), store counts from the Department of
Revenue's current licensed-dispensary list, matched to county via each city's county of location (99\%
match rate). Oregon and Montana show opposite patterns: Oregon's per-capita result is fragile (loses
significance under clustering, collapses excluding Malheur) while its per-store result is comparatively
robust (survives excluding Malheur with a \emph{lower} $p$-value); Montana's per-capita result is the only
one of the four states' per-capita estimates to remain significant under county-clustered SEs, while its
per-store result is a precise null. Section~\ref{sec:montana_mechanism} discusses why this reversal is
consistent with the license-cap mechanism proposed for Colorado and Washington.
\end{minipage}
\end{table}


Table~\ref{tab:or_mt_standalone} reports both states' revenue-per-capita and revenue-per-store results on
their own. Oregon's revenue-per-capita coefficient is large but fragile -- it does not survive
county-clustered SEs, and collapses toward zero once Malheur County is excluded, unlike its
revenue-per-store counterpart. Montana's pattern runs the other way entirely.

\subsubsection{Montana's reversed pattern: a within-design test of the license-cap mechanism}
\label{sec:montana_mechanism}

Montana's revenue-per-capita coefficient ($b=+0.654$, county-clustered $p=0.036$) is the only one of the
four states' per-capita estimates in this paper to remain significant under clustering -- yet its
revenue-per-store coefficient is a precise null ($b=+0.012$, $p=0.958$). This is the exact opposite of the
pattern this paper's central result rests on for Colorado, Washington, and Oregon, where per-capita revenue
is null and per-store revenue is where the border premium shows up. We think this reversal is not noise
but a genuine test of the license-cap explanation proposed in Section~\ref{sec:revenue_perstore}: that
explanation predicts captive border demand should surface as higher revenue per \emph{existing} store only
when the number of stores cannot freely adjust to meet it. Colorado's original licensing structure and
Washington's explicit statewide cap of 334 retail licenses (Section~\ref{sec:lottery}) both constrain
entry; Montana, by contrast, places no statewide numerical limit on the number of dispensary locations a
license holder may open, though some individual cities and counties impose their own local caps
\citep{MTCannabisLicensing}. Where entry is unconstrained, the same captive border demand that raises
revenue per store in Colorado, Washington, and Oregon should instead be absorbed by \emph{more stores}
opening in border counties -- raising revenue per capita (more total business is done) without raising
revenue per store (each individual outlet does not do disproportionately more business, because the market
was free to add outlets to meet demand). Montana's data show exactly this pattern. We offer this as a
within-design robustness check on the license-cap mechanism itself, not as an independent fourth finding:
it is the same underlying story -- captive cross-border demand -- expressed through a different margin of
adjustment because the regulatory environment differs.

\subsection{Robustness: sensitivity to the 70-mile threshold}
\label{sec:threshold}

The 70-mile cutoff used throughout this paper is a modeling choice, not a fact given by the data, and
Section~\ref{sec:data} shows it is close enough to binding for Nevada -- whose Lincoln and Elko Counties
sit 79.7 and 90.8 miles from the nearest non-legal border -- that the choice of threshold plausibly matters.
We rerun the main specification for each outcome with usable Nevada coverage at both 70 and 90 miles.

\begin{table}[htbp]
\centering
\caption{Robustness: 70-Mile vs.\ 90-Mile Border Threshold}
\label{tab:threshold_sensitivity}
\begin{tabular}{lcccccc}
\toprule
Outcome & Threshold & Coefficient & SE (HC1) & $p$ (HC1) & SE (county-cl.) & $p$ (county-cl.) \\
\midrule
Land value (23 states) & 70mi & $-0.005$ & $0.020$ & $0.789$ & $0.032$ & $0.869$ \\
Land value (23 states) & 90mi & $+0.016$ & $0.021$ & $0.448$ & $0.035$ & $0.644$ \\
\midrule
Revenue per capita (CO+WA+OR+MT) & 70mi & $+0.215$ & $0.083$ & $0.009$ & $0.193$ & $0.264$ \\
Revenue per capita (CO+WA+OR+MT) & 90mi & $+0.132$ & $0.067$ & $0.050$ & $0.152$ & $0.386$ \\
\midrule
Revenue per store (CO+WA+OR) & 70mi & $+0.429$ & $0.074$ & $<0.001$ & $0.157$ & $0.006$ \\
Revenue per store (CO+WA+OR) & 90mi & $+0.248$ & $0.068$ & $<0.001$ & $0.164$ & $0.129$ \\
\bottomrule
\end{tabular}
\begin{minipage}{0.95\textwidth}
\vspace{0.3em}
\footnotesize
\textit{Notes:} Motivated by a direct check of Nevada's computed distances (Section~\ref{sec:data}):
Lincoln County, NV sits 79.7 miles from Utah and Elko County, NV sits 90.8 miles from Idaho, both just
outside the paper's 70-mile cutoff, so Nevada's ``zero border counties'' status is a near-miss rather than
a wide margin. Widening the threshold to 90 miles pulls Nevada's Lincoln County into the treated group for
land value (Elko lacks a matched land-value record in the relevant years) and adds interior-to-border
reclassifications within the existing revenue states. Land value's null is unaffected by the threshold
choice -- both point estimates are small and both are statistically indistinguishable from zero, so the
wider net does not manufacture a land-value effect that a stricter cutoff was hiding. The two revenue
outcomes tell different stories: revenue per capita's already-fragile positive coefficient weakens further
at 90 miles and remains insignificant under clustering at both thresholds, while revenue per store's
robust 70-mile result \emph{loses} significance under county-clustered SEs at 90 miles ($p=0.006\to
p=0.129$), indicating the per-store finding, while robust to which specific counties or states are
included at a fixed 70-mile cutoff (Sections~\ref{sec:revenue_perstore} and
\ref{tab:pooled_perstore_or}), is more sensitive to the cutoff distance itself than to sample composition.
We report this openly rather than treat 70 miles as a free parameter to be tuned toward significance.
\end{minipage}
\end{table}


Land value is the one outcome where Nevada actually has data, since NASS's Census of Agriculture is a
federal program covering every county regardless of state; widening the threshold pulls Lincoln County
into the treated group there. The land-value null is unaffected: both the 70-mile and 90-mile point
estimates are small and statistically indistinguishable from zero, so a wider net does not reveal an
effect the 70-mile cutoff was suppressing. Nevada has no usable revenue data at any threshold
(Section~\ref{sec:data_revenue}), so the revenue checks in Table~\ref{tab:threshold_sensitivity} instead
test whether widening the net within Colorado, Washington, Oregon, and Montana changes those results. The
already-fragile revenue-per-capita coefficient weakens further at 90 miles. More notably, the
revenue-per-store result -- robust to excluding individual outlier counties and robust to adding an
entirely new state (Section~\ref{sec:or_mt}) -- loses significance under county-clustered SEs at 90 miles
($p=0.006\to p=0.129$). We read this as a genuine limitation rather than a reason to abandon the 70-mile
specification: the result is sensitive to how far the captive-demand radius is assumed to extend, even
though it is not sensitive to which counties or states populate a fixed radius. We report this openly, and
note that resolving it would require external evidence on how far cannabis consumers actually travel to
buy legally -- data this paper does not have -- rather than a further search for the threshold that
maximizes significance.

\subsection{A second placebo check: alcohol sales}
\label{sec:placebo_alcohol}

The general-retail placebo in Section~\ref{sec:placebo} tests whether border counties simply transact more
retail business of any kind; it does not test whether they specifically transact more business in
\emph{other regulated, excise-taxed goods with their own cross-border shopping incentive}, which is a
closer analogue to cannabis than ordinary retail and a harder placebo to pass. Alcohol is the natural
candidate: it is similarly age-restricted, similarly taxed, and has its own long-documented literature on
consumers crossing jurisdictional borders for lower prices or different availability. Colorado and
Washington, the two states behind the main revenue-per-store result, turn out to have no usable
county-level alcohol data at all (Table~\ref{tab:placebo_alcohol}'s notes detail why), so this check instead
uses Oregon and Montana -- the two states added to the design in Section~\ref{sec:or_mt} -- both of which
happen to be control states for spirits distribution with a regulator that already supplies this paper's
cannabis figures.

\begin{table}[htbp]
\centering
\caption{Placebo Check: Alcohol Sales, Oregon and Montana}
\label{tab:placebo_alcohol}
\begin{tabular}{lcccc}
\toprule
Specification & Coefficient & SE (HC1) & SE (county-cl.) & $p$-value (county-cl.) \\
\midrule
\multicolumn{5}{l}{\textit{OR+MT pooled -- outcome: log(alcohol revenue per capita), 2018--2025, $n=498$}} \\
No controls & $-0.133$ & $0.080$ & $0.194$ & $0.494$ \\
+ log(pop) + year FE & $+0.076$ & $0.081$ & $0.189$ & $0.689$ \\
+ log(pop) + year FE + state FE & $+0.209$ & $0.070$ & $0.169$ & $0.216$ \\
\midrule
\multicolumn{5}{l}{\textit{By state -- outcome: log(alcohol revenue per capita), + log(pop) + year FE}} \\
Montana (12 border counties, $n=282$) & $+0.271$ & $0.089$ & $0.217$ & $0.211$ \\
Oregon (3 border counties, $n=216$) & $-0.026$ & $0.072$ & $0.175$ & $0.883$ \\
\midrule
\multicolumn{5}{l}{\textit{Oregon only -- outcome: log(alcohol revenue per store), 2020--2025, $n=216$}} \\
No controls & $-0.848$ & $0.098$ & $0.240$ & $0.0004$ \\
+ log(pop) + year FE & $-0.245$ & $0.065$ & $0.159$ & $0.123$ \\
\bottomrule
\end{tabular}
\begin{minipage}{0.95\textwidth}
\vspace{0.3em}
\footnotesize
\textit{Notes:} Alcohol is a sharper placebo candidate than the general-retail check in
Table~\ref{tab:placebo_general_retail}: like cannabis, it is a regulated, excise-taxed good with its own
well-documented cross-border tax-arbitrage literature, so a null here is harder to attribute to simply not
being a comparable enough product category. Oregon: OLCC's own monthly, store-level spirits sales series
(the same agency that regulates cannabis retail; a control state for spirits distribution), native county
field, all 36 counties, 2019--2026 raw data aggregated to complete calendar years. Montana: also a control
state for spirits (private ``agency liquor stores'' under contract, same Alcoholic Beverage Control
Division that supplies this paper's cannabis figures); no bulk API exists, so this uses the county-level
schedule published in each year's Enterprise Fund Report of Operations (revenue.mt.gov), for the 6 fiscal
years with a machine-readable table (FY2018--2020, FY2022--2024; FY2021's schedule is a scanned image).
Parsed county totals were checked directly against each report's own printed statewide total and match
exactly in 5 of 6 years; FY2020 is 0.6\% below the report's own total; we could not attribute this to any
single county's figure and treat it as ordinary rounding in the source document rather than a parsing
error. Colorado and Washington have no usable county-level alcohol data of any kind: Colorado's liquor
excise tax series (data.colorado.gov) is statewide-only, and its general sales-tax filings (used in
Table~\ref{tab:placebo_general_retail}) do not break out a liquor-specific industry code; Washington's own
``Liquor Sales, Revenue, and Distributions'' dashboard is a non-functioning embed (confirmed directly), and
its Department of Revenue's spirits-tax statistics, like its beer and wine sales reports, are published
statewide only. \textbf{The result is mixed, not a clean null.} Oregon alone reproduces the general-retail
placebo's pattern closely -- both alcohol outcomes are null-to-negative in border counties, and if anything
alcohol revenue per store is significantly \emph{lower} near the border under HC1 (though not under
county-clustered SEs, with only 3 border-county clusters), the opposite sign from the cannabis
revenue-per-store premium. Montana shows a positive, marginally-sized border/interior gap in alcohol
revenue per capita that is significant under HC1 but not under county-clustered SEs. We do not read
Montana's result as undermining the cannabis findings -- it is a different state, product, and regulatory
structure, and the coefficient does not survive the paper's preferred (clustered) inference standard -- but
we report it rather than omit it: unlike the Colorado general-retail check, this placebo does not return a
uniformly precise null across every state and specification tried.
\end{minipage}
\end{table}


Unlike the general-retail placebo, this one does not return a clean, uniform null.
Table~\ref{tab:placebo_alcohol} shows Oregon largely replicating the general-retail pattern: alcohol
revenue per capita is a precise null in Oregon's border counties, and alcohol revenue per store is, if
anything, significantly \emph{lower} near the border under HC1 standard errors -- the opposite sign from
the cannabis revenue-per-store premium -- though this loses significance under county-clustered SEs, where
only 3 border-county clusters make the clustered estimate imprecise. Montana tells a different story:
alcohol revenue per capita is higher in its border counties, significant under HC1 ($p=0.002$) though not
under the paper's preferred county-clustered inference ($p=0.211$). We do not think this overturns the
cannabis result -- Montana's alcohol regulatory structure (state-contracted agency stores) differs from its
cannabis market, the coefficient does not survive clustering, and Oregon's own alcohol placebo still
supports the cannabis-specific interpretation -- but we report it directly rather than describing this as
another clean pass. Taken together with the general-retail check, we read the placebo evidence as
supportive of, but not uniformly conclusive for, a cannabis-specific rather than a generic border-county
interpretation of the revenue-per-store premium.

\subsection{A causal check: does the border premium switch off when the neighbor legalizes, and did sellers act on it?}
\label{sec:switchoff}

Every result so far compares border and interior counties within a fixed geographic snapshot: border
status never changes over the sample period, so none of it can rule out some other, time-invariant
difference between border and interior counties as the true cause of the revenue-per-store premium. This
section adds three checks that do not have that limitation, motivated directly by
\citet{HansenMillerWeber2020}, who show that Washington border-with-Oregon retailers' sales fell 36 percent
immediately after Oregon's own retail market opened October 1, 2015 -- i.e., that the captive-demand
advantage is not just correlated with the border, but causally switches off once the neighboring state's
own consumers no longer need to cross the line.

\subsubsection{Event-study replication (WA-OR) and a new extension (CO-NM)}

We first ask whether our own revenue-per-store construction, built independently and at annual rather than
monthly frequency, recovers a version of \citeauthor{HansenMillerWeber2020}'s WA-OR result, and then apply
the same switch-off logic to an event they did not study: New Mexico's own retail market opened April 1,
2022, which should erode the premium of the six Colorado counties bordering New Mexico (Archuleta, Conejos,
Costilla, La Plata, Las Animas, Montezuma) relative to Colorado's interior and relative to Colorado counties
that still border a permanently non-legal state (Utah, Wyoming, Oklahoma, Nebraska, Kansas).

\begin{table}[htbp]
\centering
\caption{Neighbor-State Legalization as a Switch-Off: Difference-in-Differences Estimates}
\label{tab:switch_eventstudy}
\begin{tabular}{lccccc}
\toprule
Setting & Event & Coef.\ (log rev./store) & \% change & SE (county-cl.) & $p$ (county-cl.) \\
\midrule
\citet{HansenMillerWeber2020}: WA border w/ OR & OR opens 2015-10-01 & --- & $-36\%$ & --- & --- \\
\quad (monthly diff-in-discontinuities, their data) & & & & & \\
\midrule
This paper: WA border w/ OR, annual FY panel & & & & & \\
\quad post = FY2016+ (incl.\ transition yr) & OR opens 2015-10-01 & $+0.306$ & $+35.7\%$ & $0.479$ & $0.524$ \\
\quad post = FY2017+ (FY2016 dropped) & OR opens 2015-10-01 & $+0.310$ & $+36.3\%$ & $0.492$ & $0.529$ \\
\midrule
This paper: CO border w/ NM, annual panel & & & & & \\
\quad post = 2023+ (2022 dropped) & NM opens 2022-04-01 & $-0.487$ & $-38.6\%$ & $0.123$ & $0.0001$ \\
\bottomrule
\end{tabular}
\begin{minipage}{0.95\textwidth}
\vspace{0.3em}
\footnotesize
\textit{Notes:} All rows use this paper's own log(revenue per store) construction and the project's
standard 70-mile border threshold, applied to a neighbor-state legalization date rather than to the
static distance measure used elsewhere in this paper. WA-OR: 10 border counties (Asotin, Benton, Clark,
Cowlitz, Klickitat, Lewis, Pacific, Thurston, Wahkiakum, Walla Walla) vs.\ 3 Idaho-border control counties
(Spokane, Stevens, Whitman) and 16 interior counties, $n=314$ (main) / $286$ (FY2016 dropped), controlling
log(population) and fiscal-year fixed effects. WA's revenue data is annual (fiscal year, July--June) rather
than monthly, so FY2016 (Jul 2015--Jun 2016) straddles Oregon's Oct-1-2015 opening and is a genuine
transition year, not a clean pre- or post-period; both the including- and excluding-FY2016 specifications
are reported and agree closely with each other, but neither is statistically distinguishable from zero.
CO-NM: 6 NM-border counties (Archuleta, Conejos, Costilla, La Plata, Las Animas, Montezuma) vs.\ 31 interior
counties (a parallel estimate against the 5 counties still bordering a permanently non-legal state --
Utah, Wyoming, Oklahoma, Nebraska, Kansas -- is reported in the text), $n=352$, 2022 dropped as a transition
year since New Mexico's retail market opened April 1 of that year, controlling log(population) and
year fixed effects. Figures~\ref{fig:co_nm_eventstudy} and \ref{fig:wa_or_eventstudy} show the full
event-year coefficient paths underlying these single-coefficient summaries.
\end{minipage}
\end{table}


The two settings diverge sharply. The Colorado-New Mexico extension works cleanly: dropping 2022 as a
transition year, NM-border counties' revenue per store falls $38.6\%$ relative to Colorado's interior after
New Mexico opens ($p=0.0001$, county-clustered), a magnitude close to \citeauthor{HansenMillerWeber2020}'s
own $36\%$ despite being identified from an entirely different border, state, and neighboring-state
legalization event. The event-study path in Figure~\ref{fig:co_nm_eventstudy} shows this is not a single
coefficient's artifact: relative to the reference year (2021), the NM-border premium is positive and
often significant throughout the eight pre-years, then declines monotonically from event-year 0 onward to a
precise near-zero by event-year 3; compared against the still-illegal-border control group instead of the
interior, the gap swings from parity in the pre-period to a large, significant, and widening deficit by
event-year 3 ($-1.08$ log points), consistent with NM-border counties falling behind counties that keep
their captive-demand advantage. The Washington-Oregon replication does not confirm this pattern: the point
estimate is wrong-signed (positive, i.e.\ OR-border counties' revenue per store rises rather than falls
after Oregon opens) and statistically indistinguishable from zero in both specifications
($p=0.52$--$0.53$, county-clustered), and Figure~\ref{fig:wa_or_eventstudy} shows the OR-border coefficient
climbing steadily across the full decade rather than dropping at the 2015 switch. We do not read this as
evidence against \citeauthor{HansenMillerWeber2020}'s finding, which rests on monthly administrative data
and a difference-in-discontinuities design built for exactly this transition; we read it as a limitation of
this paper's coarser annual, current-dispensary-count construction, which cannot separate a single sharp
month-level shock from a decade of differential entry and market growth at Washington's large,
population-heavy OR-border counties (Clark County alone, across the river from Portland, has 521,000
residents -- an order of magnitude larger than the true small border towns in \citeauthor{HansenMillerWeber2020}'s
setting). Excluding Washington's two largest OR-border counties (Clark and Thurston) does not change the
sign or significance. The Colorado-New Mexico setting avoids this confound because it has eight pre-event
years of data rather than one, and because its border counties are small, rural, and geographically
comparable to the interior comparison group. Taken together, we treat the CO-NM result as this paper's
strongest causal evidence for the captive-demand mechanism, and the WA-OR replication as an honest data-
resolution limitation rather than a genuine disconfirmation.

\begin{figure}[htbp]
\centering
\includegraphics[width=0.95\textwidth]{figures/fig11_co_nm_eventstudy.png}
\caption{Event-study coefficients, log(revenue per store) for Colorado counties bordering New Mexico,
relative to 2021, shown against both an interior-county and an other-nonlegal-border-county comparison
group. Dashed vertical line marks New Mexico's 2022-04-01 retail opening. Shaded bands are 95\% confidence
intervals using county-clustered standard errors.}
\label{fig:co_nm_eventstudy}
\end{figure}

\begin{figure}[htbp]
\centering
\includegraphics[width=0.85\textwidth]{figures/fig12_wa_or_eventstudy.png}
\caption{Event-year coefficients, log(revenue per store) for Washington counties bordering Oregon,
relative to FY2015 (the only fully pre-Oregon-legal year in this annual panel). Dashed vertical line marks
Oregon's 2015-10-01 retail opening, which falls inside FY2016.}
\label{fig:wa_or_eventstudy}
\end{figure}

\subsubsection{Did sellers race for the opportunity, or only respond to it after the fact?}

The revenue-per-store result says demand responds to the border; it does not say whether sellers perceived
that opportunity early and moved to capture it, or whether entry simply happened wherever it happened and
demand sorted itself out afterward. We test entry \emph{timing} directly using each dispensary's own
open date (Section~\ref{sec:data}'s geocoded dispensary panel), which the static entry-density check never
used, via Cox proportional-hazards models stratified by state.

\begin{table}[htbp]
\centering
\caption{Entry Timing: Cox Proportional-Hazards Models}
\label{tab:entry_hazard}
\begin{tabular}{lccccc}
\toprule
Outcome & Border measure & Hazard ratio & Coef. & SE & $p$ \\
\midrule
Time to 1st dispensary & border $\le$70mi (binary) & $1.138$ & $+0.130$ & $0.135$ & $0.337$ \\
(state-launch clock, $n=479$, 346 events) & log(dist.\ to non-legal border) & $0.974$ & $-0.027$ & $0.107$ & $0.803$ \\
\midrule
Time from 1st to 2nd dispensary & border $\le$70mi (binary) & $1.259$ & $+0.230$ & $0.164$ & $0.159$ \\
($n=346$, 271 events) & log(dist.\ to non-legal border) & $0.696$ & $-0.362$ & $0.124$ & $0.003$ \\
\bottomrule
\end{tabular}
\begin{minipage}{0.95\textwidth}
\vspace{0.3em}
\footnotesize
\textit{Notes:} All models stratified by state (state fixed effects) and control for log(population); the
9-state, 479-county entry-density universe and county-level open-date data described in
Section~\ref{sec:data} and used in the companion entry-density analysis. ``Time to 1st dispensary'' is
measured in days from a county's own state's retail-launch date; counties with no dispensary as of the
2026-08-05 data cutoff are right-censored at that date rather than dropped. ``Time from 1st to 2nd
dispensary'' restricts to the 346 counties that ever received a first entrant and is right-censored for
counties still at exactly one dispensary. A hazard ratio above 1 means faster entry (shorter survival time)
for counties closer to a non-legal border. The 546 dispensary records with a missing open date (including
New York's $\sim$485 license-valid-but-undated records, included in the static count analysis via lat/lon
alone) cannot be ordered in event time and are excluded here.
\end{minipage}
\end{table}


The evidence is mixed and, we think, informative precisely because of the contrast. There is no evidence
that border counties get their \emph{first} dispensary faster: both the binary and continuous border
measures are small in magnitude and statistically indistinguishable from no effect ($p=0.34$ and $p=0.80$),
consistent with the existing entry-density null -- first movers do not appear to race for border counties
specifically. But once a border county has a first entrant, a \emph{second} one follows significantly
faster the closer that county is to the border: each log-mile increase in distance to the non-legal border
lowers the hazard of a second entrant by roughly 30 percent ($p=0.003$), though the corresponding binary
measure is directionally consistent but not significant ($p=0.16$). Figure~\ref{fig:entry_hazard_km} shows
this as Kaplan-Meier survival curves: the border and interior curves for time-to-first-entrant are close
together throughout, while the second-entrant curves separate visibly, with border counties' curve falling
faster. We read this as evidence that sellers do not treat border counties as an obviously better
opportunity from the outset -- consistent with the paper's other null entry-timing and entry-density
results -- but that observed demand at an existing border store attracts a faster competitive response than
the same signal would in an interior county, a narrower and more specific version of ``competing for the
market'' than a simple race-to-enter story would predict.

\begin{figure}[htbp]
\centering
\includegraphics[width=0.95\textwidth]{figures/fig13_entry_hazard_km.png}
\caption{Kaplan-Meier survival curves, border ($\le$70mi) vs.\ interior counties, for time to first
dispensary (left) and time from first to second dispensary (right).}
\label{fig:entry_hazard_km}
\end{figure}

\subsubsection{Do dispensaries physically locate closer to the line?}

Finally, we ask whether sellers' \emph{locations within} a border county reveal the same strategic response
that their revenue does, using each dispensary's exact coordinates rather than county centroids for the
first time in this paper. For every border county, we compute the ratio of its dispensaries' mean distance
to the nearest non-legal state's polygon boundary to the county's own centroid-to-boundary distance --
a ratio below 1 means stores sit closer to the line than the county's geographic center does -- and compare
it to the same ratio computed for interior counties as a placebo.

\begin{table}[htbp]
\centering
\caption{Store-Level Bunching Toward the Non-Legal State Line}
\label{tab:border_bunching}
\begin{tabular}{lccc}
\toprule
Group & Counties & Stores & Mean store-dist.\ / county-centroid-dist.\ ratio \\
\midrule
Border counties ($\le$70mi) & 133 & 1{,}542 & $0.604$ (median $0.681$) \\
Interior counties (placebo) & 215 & 4{,}162 & $0.867$ (median $0.878$) \\
\midrule
\multicolumn{4}{l}{Welch $t$-test, border vs.\ interior ratio: $t=-10.88$, $p=6.8\times10^{-21}$} \\
\bottomrule
\end{tabular}
\begin{minipage}{0.95\textwidth}
\vspace{0.3em}
\footnotesize
\textit{Notes:} For every geocoded, currently-open dispensary (same definition as
Section~\ref{sec:data}'s entry-density count), we compute the point-to-polygon distance in miles from the
store's exact coordinates to the dissolved state polygon of its county's nearest non-legal state
(projected to EPSG:5070, CONUS Albers equal-area), and compare it to that same county's own centroid-to-
polygon distance (the measure used everywhere else in this paper). The ratio is the county's mean
store-to-line distance divided by its centroid-to-line distance; a ratio below 1 means stores sit closer to
the line, on average, than the county's own geographic center does. A simpler ``share of stores on the
border-facing side of the county centroid'' statistic was also computed but is not reported as a finding:
it is $94.6\%$ in border counties versus $91.8\%$ in interior counties -- both far above the 50\% a random
placement would imply, which shows this particular statistic is dominated by a generic, border-irrelevant
fact (incorporated towns are essentially never sited at a county's literal geometric centroid) rather than
anything specific to cross-border demand. The ratio measure controls for this by asking how much closer,
relative to the county's own scale, and it does cleanly separate the two groups: interior counties'
dispensaries sit only marginally closer to their (typically very distant and economically irrelevant)
nearest non-legal state than the county center does, while border counties' dispensaries sit
substantially closer to the actual, economically relevant line.
\end{minipage}
\end{table}


Border-county dispensaries sit, on average, at only $60\%$ of their county centroid's distance to the line,
versus $87\%$ for interior counties' dispensaries relative to their own (typically distant and economically
irrelevant) nearest non-legal state -- a large, highly significant difference ($p=6.8\times10^{-21}$). A
cruder version of this exercise, the raw share of a county's stores located on the border-facing half of
the county, is not useful on its own: it is $94.6\%$ in border counties but nearly as large, $91.8\%$, in
interior counties, which shows it is dominated by a generic fact about where towns sit within a county
rather than anything specific to the border. The ratio measure, benchmarked against each county's own
scale, is not subject to that artifact and does show real bunching: Figure~\ref{fig:border_bunching}'s left
panel plots the raw distribution of border-county store-to-line distances directly, showing a concentration
well inside the 70-mile threshold rather than a uniform spread, while the right panel shows the ratio
distributions for the two groups are clearly separated rather than overlapping. This is, to our knowledge,
the first direct, store-level (rather than county-aggregate) evidence in this research program that sellers
physically position themselves to intercept cross-border demand, complementing the revenue-per-store result
with a locational choice that is harder to attribute to anything other than proximity to the captive market
itself.

\begin{figure}[htbp]
\centering
\includegraphics[width=0.95\textwidth]{figures/fig14_border_bunching.png}
\caption{Left: histogram of individual dispensaries' distance to the nearest non-legal state's polygon
boundary, border counties only. Right: county-level distribution of the store-distance/centroid-distance
ratio, border vs.\ interior counties.}
\label{fig:border_bunching}
\end{figure}

\section{Discussion}

The four outcomes tested here do not tell a single, uniform story, and we think the differences across
them are more informative than a simple tally of nulls versus significant results would suggest. Revenue
per capita is a fragile-to-null result overall: it is a clean, well-powered null in Colorado and
Washington, weak and outlier-dependent in Oregon, and -- uniquely -- the one significant, clustering-robust
positive result in Montana, whose unlicensed-cap entry regime lets border demand show up through more
stores rather than more revenue per store (Section~\ref{sec:montana_mechanism}). Revenue per store is the
outcome with the most consistent result across states: significant and positive in Colorado, Washington,
and Oregon independently, robust to the single-county outlier in each of the two states where that check
is possible, strengthened rather than weakened by pooling a third, independently sourced state
(Section~\ref{sec:or_mt}), and -- unlike every other border-competition result in this paper or its
companion analyses -- one that \emph{strengthens} rather than fades across a state's market lifecycle in
Colorado and Washington. It is, however, sensitive to the 70-mile threshold itself
(Section~\ref{sec:threshold}), which we report as a genuine limitation rather than smooth over. The two
mechanism checks in Section~\ref{sec:mechanism} add real but unequal support: a population-weighted
captive-market-size measure survives county-clustered inference, while a simpler distance dose-response
does not, so we would characterize the mechanism evidence as meaningful but short of definitive -- direct
customer-origin data, which we do not have, would settle the question more cleanly than any proxy built
from geography alone. Read together, these revenue results suggest that the captive-demand hypothesis may
be substantively correct at the level of individual stores, but structurally invisible at the
county-aggregate level specifically where the number of stores is constrained by a policy choice (a
licensing cap) rather than left to adjust as a market outcome -- Montana's reversed pattern is the clearest
evidence for this reading, since it is the one state in the sample without such a cap, and the one state
where the per-capita/per-store roles of significance and nullity swap places.

The land-value result is null in its primary specification but is not simple: it depends on functional
form (binary threshold versus continuous distance disagree in sign), and its event study by market age
reveals a real, significant, growing negative gap several years into a state's legal market -- but only
after also revealing that the pre-launch period is itself significantly volatile in both directions, which
is precisely the kind of parallel-trends complication this research program reports transparently rather
than omits. We do not think the post-launch decline can be read as a confident causal estimate on its own;
it runs counter to the demand-absorption hypothesis, and unlike the pre-launch swings it is monotonic and
growing rather than sign-switching, but we stop short of asserting it is definitely a real effect of legal
market access rather than a continuation of whatever generates the pre-launch volatility. The
crop-substitution result is not really a null at all in the
informative sense -- it is a design that could not be tested with adequate power in the one state where we
have cultivation data, because that state's actual non-legal border is geographically distant from its
population and from where cultivation has historically concentrated.

Extending the revenue-per-store design to Oregon and Montana (Section~\ref{sec:or_mt}) was itself this kind
of test, and it held up: the mechanism identified in Section~\ref{sec:revenue_perstore} predicts that any
state with a binding license cap and genuine non-legal-border variation should show the same pattern, and
Oregon does, while Montana's absence of a statewide cap produces the theory's own predicted reversal rather
than a failure to replicate. A natural further step, beyond the scope of this paper, would be extending the
design again to additional states with binding license caps and county-level regulator data -- Michigan and
New Jersey both border non-legal states and were not reachable within this project's data-collection effort
(Section~\ref{sec:data_revenue}) but may become tractable as their own open-data infrastructure matures.
Finding cultivation data for a state with a more extensive non-legal border -- Illinois or Ohio, both of
which have substantial county-level border variation at the 70-mile threshold -- would also be considerably
better candidates than Oregon for the crop-substitution question, if comparable license data can be
obtained.

Section~\ref{sec:switchoff}'s three checks strengthen the causal case for revenue per store specifically,
though unevenly. The Colorado-New Mexico switch-off is, in our view, the single most convincing piece of
evidence in this paper: it does not rely on the assumption that border and interior counties are otherwise
comparable, only on a within-border-county comparison before and after a specific, externally-timed policy
change, and it recovers a magnitude close to \citet{HansenMillerWeber2020}'s independently-estimated effect
at a different border entirely. The Washington-Oregon replication of that same design failed on our data,
which we attribute to this paper's coarser annual, current-count revenue construction rather than to a
genuine absence of the mechanism there, since \citeauthor{HansenMillerWeber2020}'s own monthly data already
established it directly. The entry-timing and store-location checks add a second dimension the revenue
results alone cannot: sellers do not appear to race to open a border county's \emph{first} dispensary
faster than an interior county's, but once a border county has one, competitors follow measurably faster,
and dispensaries within border counties sit measurably closer to the state line than the county's own
geography would predict. Read together with the revenue-per-store result, this is a more complete picture
than the revenue evidence alone could offer: not merely a demand-side association between border proximity
and store revenue, but direct timing and locational evidence that sellers respond to, rather than are
merely correlated with, the underlying captive-demand opportunity.

\section{Conclusion}

Using the entry-game/border-competition treatment definition -- counties within 70 miles of a non-legal
state, inside an already-legal state, compared against that same state's interior counties -- we find a
mixed picture rather than a uniform null. Cannabis revenue per capita, farmland value, and crop
substitution show no robust evidence of a captive-demand effect in Colorado, Washington, and Oregon, though
for different reasons: fragile-to-null revenue-per-capita results; a null accompanied by real specification
sensitivity and a market-age pattern whose interpretation is complicated by significant pre-launch
volatility for farmland value; and an underpowered, geographically compromised design for crop substitution
in the one state we could test. Montana is the exception -- its revenue-per-capita result is the one
county-aggregate measure in this paper that is both positive and significant under county-clustered
inference, which Section~\ref{sec:montana_mechanism} argues is itself informative: Montana is also the one
state in the sample without a statewide cap on dispensary licenses, so captive border demand there can be
absorbed by new store entry rather than by higher revenue per existing store. Revenue \emph{per dispensary
store}, in the three capped states we examined, is significantly and robustly higher in border counties,
strengthens when pooled across Colorado, Washington, and Oregon rather than weakening, holds up in an event
study that strengthens rather than fades with market age in Colorado and Washington, and is supported --
though not conclusively proven -- by a population-weighted mechanism check that survives county-clustered
inference where a simpler distance measure does not; it is, however, sensitive to the exact 70-mile
threshold, a limitation we report rather than obscure. We read the combined per-store and Montana
per-capita results as the paper's central finding: the captive cross-border demand hypothesis that
motivated this entire research thread is not rejected once the outcome is measured at the right level of
aggregation given each state's own licensing regime, even though every county-aggregate measure tried
elsewhere in this project under a binding license cap -- entry timing, entry density, and revenue per
capita alike -- has come back null. We hold this finding with appropriate caution rather than as a settled
result: the mechanism evidence is suggestive rather than direct, and confirming it properly would require
customer-origin data this paper does not have.

\newpage
\bibliographystyle{plainnat}
\bibliography{references_substitution}

\end{document}
